\subsection{H8: the exceptional divisor lattice, with its dimension-one exception}

On 57:7 Grothendieck asks that $D\to\mathrm{NS}(X_0)$ be injective; on 57:8
he writes ``I expect that this is always the case'' and proposes intersection
with curves on the exceptional components. These leaves are in English, not
French, so a French original must not be invented.

Let $A$ be a complete noetherian local ring, let
$f:X\to\operatorname{Spec}A$ be proper, regular, and the closure of the
regular punctured spectrum, and let $D$ denote the free group of prime
divisors supported on the closed fibre. If the irreducible components of
$X$ have dimension at least two, then
\[
 D\longrightarrow\mathrm{NS}(X_0)
\]
is injective, even after replacing algebraic equivalence by numerical
equivalence. Indeed first suppose $A$ normal and integral. A divisor $B$
whose restriction to $X_0$ is algebraically trivial has degree zero on every
proper curve contracted by $f$: all those curves lie in $X_0$. Both $B$ and
$-B$ are therefore $f$-nef, and $f_*B=0$ because the closed point has
codimension at least two. The negativity lemma, applied to $B$ and to $-B$,
gives $B\geq0$ and $B\leq0$, whence $B=0$.
The exact source opened is Koll\'ar, \emph{Families of varieties of general
type}, Lemma~11.60, p.~436, citing Koll\'ar--Mori~3.39. Section~11.4 explicitly
treats arbitrary schemes; the characteristic-zero restriction stated at the
start of Chapter~11 is for Sections~11.1--11.3.

For general $A$, replace it by the normalization of $A_{\mathrm{red}}$.
It is finite since $A$ is excellent. The normal scheme $X$ factors through
it, and the punctured spectrum is unchanged since it was regular. The finite
normalization is a disjoint union of complete normal local schemes. Apply the
preceding argument on each component. All contracted curves remain in the
original closed fibre, so numerical triviality is preserved.

There is an indispensable qualification in dimension one. For
$A=k[[t]]$ and $X=\operatorname{Spec}A$, the divisor group is $\mathbf Z$,
whereas $\mathrm{NS}(X_0)=\mathrm{NS}(\operatorname{Spec}k)=0$. Thus the
unqualified assertion for every dimension is false. More generally, in the
setting just stated, the whole kernel is freely generated by the closed points
of the one-dimensional components of $X$: after normalization those components
are spectra of DVRs and their proper birational maps are isomorphisms; the
dimension-at-least-two components have zero kernel by negativity. Components
of dimension zero contribute no divisors. This gives the complete correction.

\subsection{H10: smoothness in characteristic zero over nonreduced bases}

The supplementary question on 57:64 reads:
\begin{quote}
``Si $G$ est un sch\'em en groupes sur $S$, $S$ \`a corps r\'esiduels de car.
$0$, et si $\underline\omega_{G/S}$ [faisceau conormal \`a la section unit\'e]
est loc. libre, alors $G$ est simple sur $S$ le long de la section unit\'e''.
\end{quote}
With the locally-finite-presentation hypothesis used throughout the discussion,
this is exactly SGA~3, VI$_B$, Proposition~1.6 (original p.~327, p.~6 of the
opened retyped expos\'e). No flatness hypothesis is needed. The proof uses
$\Omega^1_{G/S}=\pi^*\omega_{G/S}$ and EGA IV$_4$, 16.12.2 and 17.12.5.

For $X/S$ smooth projective in characteristic zero, Deligne's
\emph{Th\'eor\`eme de Lefschetz et crit\`eres de d\'eg\'en\'erescence de suites
spectrales}, Theorem~5.5(i), p.~123, gives local freeness and arbitrary
base-change compatibility of every $R^qf_*\Omega^p_{X/S}$, with no reducedness
condition on $S$. The infinitesimal Picard tangent functor is
$M\mapsto R^1f_*(\mathcal O_X\otimes f^*M)$, which, by this base-change theorem,
equals $R^1f_*\mathcal O_X\otimes M$. Consequently the conormal module of Pic
at the identity is $(R^1f_*\mathcal O_X)^\vee$ and is locally free.
The preceding SGA~3 criterion makes Pic smooth along the identity. Its neutral
part is therefore an open smooth group scheme. Its fibres are abelian varieties
(smooth proper varieties have proper reduced neutral Picard schemes), and its
relative properness follows from the normal-fibre proper-neutral-subset theorem
once that part is open and represented. Thus $\operatorname{Pic}^0_{X/S}$ is an
abelian scheme over arbitrary, possibly nonreduced, $S$.
If the geometric fibres of $X/S$ are disconnected, make the finite \'etale
Stein-factor reduction and take the corresponding product/Weil restriction.
Kleiman's Remark~5.21, pp.~48--49, provides an independent presentation of the
same deduction using exponential lifting; this is a source comparison, not a
claim of an independent research verification.

\subsection{H11: affineness after shrinking for normal families over a curve}

The question at 57:74, repeated on 57:78, is in English in the manuscript:
``one could replace `quasi-affine' by `affine'.'' It specifically proposes
finite maps between projective normal families, and birational maps between
smooth projective families. The finite-normal possibility has an affirmative
answer after the permitted shrinking; the argument below also allows proper
surjections between projective normal families. It does not assert the same
result for every nonnormal proper family.

\paragraph{Theorem H11.N.}
Let $S$ be integral noetherian of dimension at most one, and let
$f:X'\to X$ be a surjective morphism of projective $S$-schemes whose geometric
generic fibres are normal. After replacing $S$ by a nonempty open, the map
$f^*:\operatorname{Pic}_{X/S}\to\operatorname{Pic}_{X'/S}$ is affine.
In particular the finite normal-family counterexample suggested on 57:78
cannot exist after further shrinking of the base.

\paragraph{Proof.}
Shrink first so both families are flat with geometrically normal fibres,
and so both Picard functors are separated group schemes locally of finite
presentation and $u=f^*$ is quasi-affine of finite presentation, as in
SGA~6 XII Theorem~1.1 and Corollary~1.2 (also Kleiman, Theorem~4.18.5).
Write $P=\operatorname{Pic}_{X/S}$, $P'=\operatorname{Pic}_{X'/S}$,
$H=P^\tau$, $H'=P'^\tau$. These are open-and-closed group subschemes of
finite presentation. Their generic fibres are proper, since the generic
geometric fibres of $X,X'$ are normal. Generic flatness and spreading out
properness for a finitely presented morphism allow one more shrinking so
that $H,H'$ are proper and flat over $S$.

By H6 the quotients $Q=P/H$, $Q'=P'/H'$ are separated unramified schemes
locally of finite presentation over $S$. The projections $P\to Q$ and
$P'\to Q'$ are fppf torsors under $H,H'$, respectively, and hence are proper,
flat, surjective, and of finite presentation. Pullback induces $v:Q\to Q'$.

First $v$ is a monomorphism. On a geometric fibre, a line bundle whose
pullback is numerically trivial is itself numerically trivial. In the finite
case, for every integral curve $C\subset X_s$ one takes a component
$C'\subset X'_s$ finite and surjective onto $C$ and uses
\[
 \deg_{C'}f^*L=[k(C'):k(C)]\deg_C L.
\]
For a proper surjection one can use Chow's lemma and successive hyperplane
sections over $C$ to obtain such a dominating proper curve, and the same
degree formula applies. Numerical triviality is equivalent to membership
in $\operatorname{Pic}^\tau$ for a proper scheme over an algebraically closed
field. Thus $v$ is injective on geometric points. If two $T$-points of $Q$
have the same image, their values agree on every geometric fibre of $T$.
Their equalizer is open, since $Q/S$ is unramified, and contains all points
of $T$; therefore the two $T$-points are equal.

Next $v$ is quasi-compact. For any quasi-compact open $V\subset Q'$, the
scheme $P'\times_{Q'}V$ is quasi-compact because $P'\to Q'$ is proper.
Its inverse image under the quasi-compact map $u$ is quasi-compact, and maps
surjectively to $v^{-1}(V)$. Thus $v^{-1}(V)$ is quasi-compact. The morphism
$v$ is consequently separated and quasi-finite. Every local ring of $Q'$
has dimension at most one, since $Q'/S$ is unramified and $\dim S\le1$.
Zariski's Main Theorem embeds $v^{-1}(V)$, for affine $V$, as a quasi-compact
open in a finite $V$-scheme. Such a finite scheme is affine and noetherian
of dimension at most one, and every quasi-compact open in it is affine
(Stacks, Lemma~33.38.2, tag~09N9). Thus $v$ is affine.

Finally consider
\[
 P\xrightarrow{\psi} P'\times_{Q'}Q\longrightarrow P'.
\]
Both $P\to Q$ and $P'\times_{Q'}Q\to Q$ are proper, and the latter is
separated, so $\psi$ is proper. It is quasi-affine: this follows by taking
the graph over $P'$ and using that $u$ is quasi-affine and the second map
is separated. A proper quasi-affine morphism is finite. The second map
is affine as a base change of $v$, so their composite $u$ is affine.

The proof keeps the finite-type hypothesis on $H$ visible: normality alone
does not make $H$ proper over the original unshrunk base, as Theorem E4.N
shows. The smooth birational case on 57:78 is covered by the same proof.
One must not claim that the full Picard functor is unchanged by a blowup:
exceptional divisor classes enlarge the N\'eron--Severi group. The familiar
birational invariance concerns the neutral Picard variety, and pullback on
the full Picard functor is injective when $f_*\mathcal O_{X'}=\mathcal O_X$.
The argument above does not address nonnormal fibres. The following theorem
settles that branch by a different quotient, using the kernel itself.

\subsection{H11: the full one-dimensional affineness statement}

\paragraph{Theorem H11.G.}
Let $S$ be an integral noetherian scheme of dimension at most one, and let
$f:X'\to X$ be a surjective morphism of proper $S$-schemes. After replacing
$S$ by a nonempty open subscheme, the pullback
\[
 f^*:\operatorname{Pic}_{X/S}\longrightarrow\operatorname{Pic}_{X'/S}
\]
is affine. No normality or projectivity assumption is necessary.
Thus the general one-dimensional question on 57:74 and 57:78 has an
 affirmative answer, extending Theorem H11.N above.

\paragraph{Proof.}
SGA~6 XII Theorem~1.1 and Corollary~1.2,
\cite[original p.~596]{BCESGA6XII}, allow a first shrinking so that
$P=\operatorname{Pic}_{X/S}$ and $P'=\operatorname{Pic}_{X'/S}$ are separated
group schemes locally of finite presentation, and $u=f^*:P\to P'$ is
quasi-affine of finite presentation. Its kernel $K$ is a quasi-affine group
scheme of finite presentation over $S$. The generic fibre $K_\eta$ is an
affine algebraic group, by SGA~3 VI$_B$~11.11. Affineness spreads out for a
finitely presented scheme over the generic point: take an affine model of
$K_\eta$ over a nonempty open and spread out the generic isomorphism and
its inverse. After further shrinking, $K$ is therefore affine. Generic
flatness permits a final shrinking making $K/S$ flat.

The free translation action of the flat finitely presented group $K$ on
$P$ has a quotient algebraic space $Q=P/K$ (Stacks tag~06PH). The quotient
map $p:P\to Q$ is an fppf $K$-torsor, hence affine and of finite presentation
by descent. The map $u$ factors through a monomorphism $v:Q\to P'$.
This map is locally of finite presentation, by fppf descent from $u$ along
$p$, and is quasi-compact: over a quasi-compact open $V\subset P'$, the
quasi-compact scheme $u^{-1}(V)$ maps surjectively to $v^{-1}(V)$.
Thus $v$ is a quasi-finite separated morphism of finite presentation.
Stacks tag~0418 shows that it is representable by schemes, so $Q$ is a
scheme as well.

For every point $s\in S$, the morphism $v_s$ is a quasi-compact
monomorphism of group schemes locally of finite type over the field
$\kappa(s)$. SGA~3 VI$_B$, Corollary~1.4.2, original pp.325--326, says
precisely that such a monomorphism is a closed immersion. In particular
$v_s$ is finite. Now apply SGA~3 XVI, Lemma~4.2, original pp.521--522
(p.~19 of the opened retyped expos\'e): over a locally noetherian base
of dimension at most one, a quasi-finite separated morphism between
schemes locally of finite type, whose morphism on every fibre is finite,
is affine. Therefore $v$ is affine and so is $u=v\circ p$.

For clarity, the quoted lemma imposes no dimension bound on $P'$ itself.
Its proof passes to complete local bases and then a DVR. Zariski's Main
Theorem splits the source over a complete local target into a finite part
and a part missing its closed point. Fibrewise finiteness forces the latter
part to lie in the generic fibre; that generic fibre is a principal affine
open over the DVR. Both parts are affine over the target. This is the
one-dimensional input that removes the earlier normality restriction.
The statement does not assert the corresponding affineness over an
arbitrary higher-dimensional base.
